\section{Editors and IDEs}

Go is supported by a relatively broad range of integrated development environments (IDEs) and extensible text editors. Unlike some languages, however, Go does not require a heavyweight IDE: the standard Go toolchain provides a substantial set of command-line facilities for building, testing, formatting and analysing programs, whilst the official \texttt{gopls} language server supplies much of the functionality normally associated with an IDE.\\

The choice of development environment is therefore primarily a question of workflow, project complexity, and personal preference. The principal options are Visual Studio Code, GoLand, Vim/Neovim and, to a lesser extent, other editors such as Emacs and LiteIDE.\\

\subsection{Visual Studio Code}

\href{https://code.visualstudio.com/}{Visual Studio Code (VS Code)} is arguably the most widely used development environment for Go. The 2024 H2 Go Developer Survey reported that 66\% of the respondents used VS Code regularly, compared to 35\% for GoLand, making VS Code the clear leader in terms of reported usage.\\

Go support is provided principally by the official
\href{https://marketplace.visualstudio.com/items?itemName=golang.go}{Go extension}, maintained by the Go team. The extension integrates closely with \texttt{gopls}, the official Go language server, and provides IntelliSense, code navigation, diagnostics, refactoring, formatting, testing, and debugging.\\

A significant advantage of this architecture is that much of the functionality is provided by the same underlying Go tooling used by other editors. Consequently, developers can obtain a consistent development experience whilst retaining the comparatively lightweight and extensible nature of VS Code.\\

\begin{itemize}
\item \textbf{Advantages:}
\begin{itemize}
\item Free and available on Windows, macOS and Linux.
\item Excellent Go support through the official Go extension.
\item Strong integration with \texttt{gopls}, the Go language server.
\item Integrated testing and debugging.
\item Extensive ecosystem of extensions for containers, Kubernetes, Git, databases, and cloud platforms.
\item Good support for multi-language projects in which Go is used alongside other languages.
\item Relatively low barrier of entry compared to a traditional full-featured IDE.
\end{itemize}

\item \textbf{Disadvantages:}
\begin{itemize}
    \item Some functionality depends on extensions and external Go tools rather than being part of the core editor.
    \item The large extension ecosystem can result in configuration complexity.
    \item Resource consumption can increase substantially when many extensions are installed.
    \item Some advanced IDE functionality may require more manual configuration than in GoLand.
\end{itemize}

\end{itemize}

For most Go developers, VS Code is a strong default choice, particularly where a free, extensible, and multi-purpose development environment is required.\\

More information is available from the
\href{https://go.dev/doc/editors}{official Go editor documentation}
and the
\href{https://code.visualstudio.com/docs/languages/go}{VS Code Go documentation}.

\subsection{GoLand}

\href{https://www.jetbrains.com/go/}{GoLand}, developed by JetBrains, is a dedicated commercial IDE designed specifically for Go development. It is based on the IntelliJ platform and provIDE. a highly integrated development environment in which many features that require extensions or additional configuration in lighter editors are available as part of the IDE.\\

GoLand provides sophisticated code completion, navigation, refactoring, debugging, testing, profiling, and integration with Go modules. It also provides support for common technologies surrounding Go development, including version control systems, databases, HTTP clients, Docker and Kubernetes.\\

The principal distinction between GoLand and VS Code is therefore not whether either can provide the essential Go functionality, both are capable of this, but rather how much of the surrounding development workflow is integrated into the IDE itself.

\begin{itemize}
\item \textbf{Advantages:}
\begin{itemize}
\item Comprehensive, dedicated Go development environment.
\item Very strong code navigation and refactoring facilities.
\item Integrated debugging, testing and profiling.
\item Good support for Go modules and complex project structures.
\item Extensive integration with databases, containers, version control and other development tools.
\item Generally requires less extension-level configuration than VS Code.
\item Particularly well suited to large and complex commercial Go projects.
\end{itemize}

\item \textbf{Disadvantages:}
\begin{itemize}
    \item Commercial licences are required for many users.
    \item Heavier in terms of memory and system resources than a minimalist editor.
    \item Less attractive for developers who prefer a highly configurable, editor-centric workflow.
    \item Its breadth of functionality can be unnecessary for small Go programs or simple command-line utilities.
\end{itemize}

\end{itemize}

GoLand is therefore particularly attractive for professional teams working on sizeable Go codebases where sophisticated refactoring, debugging, and project-management facilities justify the additional resource requirements and licencing considerations.\\

It is worth noting that Go support is also available through the Go plugin for
\href{https://www.jetbrains.com/idea/}{IntelliJ IDEA Ultimate}; JetBrains states that the Go plugin provides the functionality available in GoLand. This can be advantageous for developers who already use the IntelliJ ecosystem for other languages.\\

More information is available from the
\href{https://www.jetbrains.com/go/}{GoLand product documentation}
and the
\href{https://www.jetbrains.com/help/idea/go-plugin.html}{IntelliJ IDEA Go documentation}.

\subsection{Vim and Neovim}

\href{https://www.vim.org/}{Vim} and
\href{https://neovim.io/}{Neovim} occupy a different position from conventional IDEs. They are highly configurable text editors that can be transformed into capable Go development environments through plugins and language-server integration.\\

The most important component is again \texttt{gopls}. Neovim provides native Language Server Protocol (LSP) facilities, whilst Vim can use a number of LSP-oriented plugins. The Go project provides specific guidance for configuring \texttt{gopls} with Vim and Neovim.\\

A typical configuration can provide completion, diagnostics, code navigation, formatting, import organisation, and other IDE-like facilities without abandoning the keyboard-driven Vim workflow.

\begin{itemize}
\item \textbf{Advantages:}
\begin{itemize}
\item Extremely fast and lightweight once configured.
\item Excellent keyboard-driven workflow.
\item Highly configurable and scriptable.
\item Particularly effective when working over SSH or on remote systems.
\item Excellent integration with command-line Go tooling.
\item No requirement to adopt a large graphical IDE.
\end{itemize}

\item \textbf{Disadvantages:}
\begin{itemize}
    \item Initial configuration can be considerably more complicated than VS Code or GoLand.
    \item IDE functionality depends on plugins and configuration choices.
    \item Debugging and project-management workflows can require additional setup.
    \item Less suitable for developers unfamiliar with modal editors.
\end{itemize}

\end{itemize}

Vim or Neovim is therefore particularly appropriate for experienced developers who value efficiency, customisation, and terminal-based workflows. It is less compelling as a first Go environment for developers who want a fully integrated graphical IDE immediately.\\

The official Go documentation contains 
\href{https://go.dev/gopls/editor/vim}{guide to using \texttt{gopls} with Vim and Neovim}.

\subsection{Emacs}

\href{https://www.gnu.org/software/emacs/}{Emacs} is another extensible editor capable of providing a sophisticated Go development environment. Go support can be built around \texttt{gopls} and the Language Server Protocol, commonly through packages such as \texttt{lsp-mode} or related LSP integrations.\\

Emacs is particularly powerful for developers who already use it as their primary environment, since Go development can be integrated with terminals, version control, documentation, project management, and other workflows within a single highly programmable environment.\\

Its main limitation is the same characteristic that makes it attractive: considerable configurability. Achieving a polished Go environment generally requires more configuration than using GoLand or the official VS Code extension.

\begin{itemize}
\item \textbf{Advantages:}
\begin{itemize}
\item Exceptionally extensible and programmable.
\item Strong LSP support through \texttt{gopls}.
\item Excellent terminal and version-control integration.
\item Well suited to developers who already use Emacs extensively.
\end{itemize}

\item \textbf{Disadvantages:}
\begin{itemize}
    \item Greater configuration overhead.
    \item Less immediately accessible to newcomers.
    \item A number of Go-specific features depend on third-party packages and configuration.
\end{itemize}

\end{itemize}

\subsection{LiteIDE and Other Editors}

\href{https://github.com/visualfc/liteide}{LiteIDE} is an open-source IDE specifically oriented towards Go. ProvIDE a simpler, more traditional IDE experience and includes features such as Go-aware editing and integration with Go tools.\\

However, LiteIDE is considerably less prominent than VS Code and GoLand in the contemporary Go ecosystem. The Go project's own documentation lists it among the available editors, but the latest Go Developer Survey indicates that usage is heavily concentrated around VS Code and GoLand, with Vim/Neovim and other editors occupying smaller niches.\\

Other editors, including Sublime Text, Zed and Kate, can also provide Go development facilities through language-server or editor integrations. These can be excellent choices for users who already prefer those environments, but they are generally less significant as dedicated Go IDE choices.

\subsection{Comparison}

Table~\ref{tab:go-ides} summarises the principal trade-offs.

\begin{table}[htbp]
\centering
\caption{Comparison of major development environments for Go}
\label{tab:go-ides}
\begin{tabular}{p{0.16\textwidth}p{0.20\textwidth}p{0.27\textwidth}p{0.27\textwidth}}
\hline
\textbf{Environment} & \textbf{Model} & \textbf{Advantages} & \textbf{Disadvantages} \
\hline
VS Code &
Extensible editor/IDE &
Free; excellent official Go extension; \texttt{gopls} integration; strong debugging and testing; extensive ecosystem &
Some functionality requires extensions; configuration can become complex; less integrated than a dedicated IDE \

GoLand &
Dedicated commercial IDE &
Excellent Go-specific tooling; advanced refactoring, debugging and profiling; strong project and database integration &
Commercial for many users; comparatively resource-intensive; potentially excessive for small projects \

Vim/Neovim &
Highly configurable editor &
Very fast; keyboard-centric; excellent terminal and remote-development workflow; highly customisable &
Steeper learning curve; significant configuration and plugin management; some IDE functionality requires manual setup \

Emacs &
Extensible editor &
Extremely programmable; strong LSP support; excellent integration with command-line workflows &
High configuration overhead; less immediately accessible; Go functionality is assembled from packages \

LiteIDE &
Dedicated open-source Go IDE &
Simple Go-focused environment; open source; relatively lightweight &
Much smaller ecosystem and user base; fewer integrations and less momentum than VS Code or GoLand \
\hline
\end{tabular}
\end{table}

\subsection{The Role of \texttt{gopls}}

An important consideration when comparing Go development environments is that the boundary between an ``IDE feature'' and the underlying Go tooling is relatively thin. The official \texttt{gopls} language server provides a common foundation for many editors.\\

\texttt{gopls} supplies facilities such as code completion, navigation, diagnostics, symbol information, refactoring, and code actions. This means that a developer moving between VS Code, Neovim, Emacs and other LSP-compatible environments can retain a substantial proportion of the same language intelligence.\\

This architecture is beneficial to the Go ecosystem because the editor choice does not necessarily imply choosing a completely different language-analysis implementation. The official VS Code extension, for example, uses \texttt{gopls} as its default backend for many IntelliSense, navigation, editing and diagnostic features.

\subsection{Recommendation}

For most developers, the choice can be reduced to three principal options.

\begin{enumerate}
\item \textbf{VS Code} is the best general-purpose choice. It is free, widely adopted, well supported by the Go team, and sufficiently powerful for both small projects and substantial production systems.

\item \textbf{GoLand} is the strongest choice when a fully integrated, dedicated IDE is preferred. It is particularly compelling for large commercial applications where advanced refactoring, debugging, profiling, and project integration are important.

\item \textbf{Vim or Neovim} is an excellent choice for experienced developers who prefer a highly customisable and terminal-centric environment, especially when remote development over SSH is common.\\

\end{enumerate}

The latest available Go Developer Survey provides useful evidence for this hierarchy: VS Code was regularly used by 66\% of respondents and GoLand by 35\%, whilst 33\% reported regularly using two or more editors. The latter figure is significant because it demonstrates that editor choice need not be exclusive; for example, a developer may use GoLand or VS Code for primary development while using Neovim or Vim for quick edits on remote systems.\\

Ultimately, the most important factor is not the IDE itself, but the quality of integration with the standard Go toolchain. A suitable environment should provide reliable access to \texttt{go}, \texttt{gofmt}, \texttt{go test}, the Go module system, \texttt{gopls} and, where required, a debugger such as Delve. On this basis, VS Code and GoLand currently offer the strongest combination of maturity, functionality, and ecosystem support, while Vim/Neovim and Emacs remain compelling alternatives for developers with established editor-centric workflows.

\subsection*{Further Information}

\begin{itemize}
\item \href{https://go.dev/doc/editors}{Go project: Editor plugins and IDEs}
\item \href{https://go.dev/wiki/IDEsAndTextEditorPlugins}{Go Wiki: Editors and IDEs}
\item \href{https://go.dev/blog/survey2024-h2-results}{Go Developer Survey 2024 H2 results}
\item \href{https://code.visualstudio.com/docs/languages/go}{VS Code: Go documentation}
\item \href{https://www.jetbrains.com/go/}{JetBrains GoLand}
\item \href{https://go.dev/gopls/}{Go language server (\texttt{gopls})}
\item \href{https://go.dev/gopls/editor/vim}{\texttt{gopls}: Vim and Neovim}
\end{itemize}